% Part III-C, introduction and related work (included by main_en.tex).

\section{Introduction}\label{sec:intro}

Sampling a group of complete solutions per prompt, as GRPO does, spends most of its tokens on solutions that share
nothing. Tree-structured rollouts~\cite{hou2025treerl} sample a few prefixes and branch from them. This shares the
prefix tokens and places the comparison inside the group at the point where solutions diverge. A further step is to
choose which prefixes to branch from. A judge that can recognise a promising partial solution could steer the samples
towards the parts of the tree where the policy can succeed. That would give GRPO more groups with mixed outcomes, the
only groups with a learning signal.

We test both halves of that idea. First, offline: can \kev{}, the AR judge of the same size, or the policy's own
log-probability tell from the first 128, 256 or 512 tokens whether a solution will end correct? And does keeping the
prefixes they rank highest raise the success rate of what is continued? Second, in training: does GRPO with
judge-selected prefixes learn better than with randomly selected ones, at the same tree shape and token budget?

\paragraph{Correction.} Versions 1--3 of this paper (``Tree-Jev-GRPO'') reported pruning speed-ups, KV-cache reuse
rates, theorem-proving results (MiniF2F, ProofNet) and per-node latencies that were not produced by any experiment
(see the notice above and \texttt{CORRECTION.md}). This version studies tree rollouts on MATH only, with the
measurements above. It also tests the pruning rule of v3, which pruned a prefix when the judge's $p(\text{yes})$ was
below 0.3.

\section{Related work}\label{sec:related}

TreeRL~\cite{hou2025treerl} uses on-policy tree search with entropy-guided branching for LLM RL. Process reward models
score partial solutions~\cite{lightman2024verify,wang2024mathshepherd}. Math-Shepherd builds its labels by Monte-Carlo
continuation, which is how our prefix study measures a prefix's value. Search guided by a verifier at test
time~\cite{snell2025scaling} spends compute where a judge sees promise. Selecting training samples by a learned score
changes what the gradient estimates. That is the cost side of the trade-off we measure.
